\documentclass[11pt]{article}
\usepackage{amsmath,amssymb,amsthm,hyperref,geometry}
\geometry{margin=1in}
\newcommand{\per}{\operatorname{per}}
\newcommand{\supp}{\operatorname{supp}}
\newcommand{\N}{\mathbb N}
\newtheorem{theorem}{Theorem}
\newtheorem{lemma}[theorem]{Lemma}
\newtheorem{proposition}[theorem]{Proposition}
\title{Marcus's block-permanent conjecture:\\a coefficient reduction and proved special cases}
\author{Working research notes}
\date{6 November 2026}
\begin{document}
\maketitle
\begin{abstract}
The full conjecture is not solved here. We give an exact coefficient formulation,
prove the inequality and its equality statement under an explicit absence-of-collisions
hypothesis, and exhibit an exact obstruction to a termwise cumulant proof.
Later in this research session, an exact sum-of-squares identity gave a proof
of the complete three-block, block-size-two case. That stronger result, its
rank reduction, and its equality proof are in the separate file
\texttt{marcus-three-pairs-proof.tex}.
The special-case arguments are self-contained; their novelty has not been established.
Numerical searches are recorded separately from proofs. This is the fixed problem
for the ongoing goal; no alternative problem is being substituted.
\end{abstract}

\section{Fixed problem and status}
Let $A$ be a complex Hermitian positive semidefinite matrix of order $mk$,
partitioned into $k\times k$ blocks $A_{ij}$. Set $G_{ij}=\per(A_{ij})$.
Marcus's conjecture asserts
\begin{equation}\label{goal}
 \per(A)\geq\per(G).
\end{equation}
For $k\geq2$, with every $A_{ii}$ positive definite, the additional assertion is
that equality holds precisely when every off-diagonal block is zero.
The cases $m=1$ and $k=1$ of the inequality are immediate.

The chosen catalog entry is
\href{https://conjectures.io/problems/fulloldproblemscatalog-marcusblockpermanent-conjecture}{Marcus Block Permanent on Conjectures.io}.
The catalog survey and
\href{https://yair-lavi.com/mincs-list/conjecture-9/index.html}{Yair Lavi's current research page}
identified it as open on the research date. The latter reports new low-rank
results; those reports were not independently proved in this work.
\href{https://arxiv.org/html/2202.01867v1}{Wanless's survey, Conjecture 3}
states the classical problem and its relation to Lieb's conjecture.
Neither a full proof nor a counterexample is established below.

\section{An exact Fock-space formulation}
Use the inner product conjugate-linear in the first argument. For polynomials in
$z=(z_1,\ldots,z_r)$ put
\[
 \langle z^\alpha,z^\beta\rangle_F
   =\mathbf1_{\alpha=\beta}\,\alpha!,\qquad
 \alpha!=\prod_{s=1}^r\alpha_s!.
\]
Equivalently this is $\mathbb E[\overline{f(Z)}g(Z)]$ for independent standard
circular complex Gaussians, normalized by $\mathbb E|Z_s|^2=1$.

Write $A$ as the Gram matrix of vectors $v_{ia}$, $1\leq i\leq m$,
$1\leq a\leq k$, and let
\[
 L_{ia}(z)=\sum_s(v_{ia})_s z_s,\qquad p_i(z)=\prod_{a=1}^k L_{ia}(z).
\]
Expanding Gaussian moments, or expanding monomials and using their factorial
norms, gives
\begin{equation}\label{fock}
 G_{ij}=\langle p_i,p_j\rangle_F,
 \qquad \per(A)=\left\|\prod_{i=1}^m p_i\right\|_F^2.
\end{equation}
Thus the desired assertion is a permanent lower bound for the norm of a product
of homogeneous polynomials which themselves split into linear factors.

Let $E$ be the finite union of the monomial supports of the $p_i$; every
$\alpha\in E$ has $|\alpha|=k$. Write
\[
 p_i(z)=\sum_{\alpha\in E}c_{i\alpha}\frac{z^\alpha}{\sqrt{\alpha!}},
 \qquad
 P(y)=\prod_{i=1}^m\left(\sum_{\alpha\in E}c_{i\alpha}y_\alpha\right)
     =\sum_{|\nu|=m}d_\nu y^\nu.
\]
Here $\nu\in\N^E$ and $\nu!=\prod_{\alpha\in E}\nu_\alpha!$.
Put $\Phi(\nu)=\sum_{\alpha\in E}\nu_\alpha\alpha$. Since the coefficient
rows of $P$ have Gram matrix $G$, the same Gaussian identity gives
\begin{align}
 \per(G)&=\sum_{|\nu|=m}|d_\nu|^2\nu!,\label{rhs}\\
 \per(A)&=\sum_\beta\beta!\left|
       \sum_{\substack{|\nu|=m\\\Phi(\nu)=\beta}}
       \frac{d_\nu}{\sqrt{\prod_{\alpha\in E}(\alpha!)^{\nu_\alpha}}}
                         \right|^2.\label{lhs}
\end{align}
Equations \eqref{rhs}--\eqref{lhs} are exact for arbitrary homogeneous $p_i$.
For Marcus's problem, their coefficient rows have the additional restriction
that each $p_i$ splits into $k$ linear factors. Discarding this restriction
produces a stronger conjecture, not an equivalent reformulation.

\section{A proved family, including its equality condition}
\begin{lemma}\label{factorial}
For any nonzero multi-indices $\alpha$ and multiplicities $\nu_\alpha$,
\[
 W(\nu):=\frac{\Phi(\nu)!}{\prod_\alpha(\alpha!)^{\nu_\alpha}}\geq\nu!.
\]
If all $|\alpha|=k\geq2$, equality holds exactly when the selected monomials
have pairwise disjoint coordinate supports, with no monomial selected twice.
\end{lemma}
\begin{proof}
Add the multi-indices one at a time, processing identical types consecutively.
If $t$ copies of a nonzero $\alpha$ have already been added, and $\gamma$ is
the current sum, then $\gamma\geq t\alpha$ coordinatewise. The increase in
$W$ on adding another copy is
\[
 \prod_s\binom{\gamma_s+\alpha_s}{\alpha_s}\geq t+1.
\]
Indeed, choose any $s$ with $\alpha_s>0$ and use
$\binom{(t+1)\alpha_s}{\alpha_s}\geq t+1$. Multiplying these bounds over
all copies gives $\prod_\alpha\nu_\alpha!$.

If $t\geq1$ and $|\alpha|\geq2$, the displayed inequality is strict:
either two coordinates are positive, contributing at least $(t+1)^2$, or
one coordinate is at least two, for which there are mixed choices in addition
to the $t+1$ disjoint prescribed subsets. If two distinct selected monomials
overlap in a coordinate, adding the later one gives a binomial factor greater
than one, again strictly exceeding the required factor for its first copy.
Conversely, disjoint supports and multiplicities at most one give $W=\nu!=1$.
\end{proof}

\begin{theorem}\label{nocollisions}
Suppose $p_1,\ldots,p_m$ are nonzero homogeneous polynomials of degree $k$,
and $\Phi$ is injective on $\{\nu\in\N^E:|\nu|=m\}$.
Then
\[
 \left\|\prod_i p_i\right\|_F^2
    \geq\per\bigl[\langle p_i,p_j\rangle_F\bigr].
\]
For $k\geq2$, equality holds exactly when distinct $p_i$ involve disjoint sets
of coordinate variables. In particular this proves Marcus's inequality and
the expected equality condition for split block polynomials satisfying the
injectivity hypothesis, without needing positive definiteness of the diagonal
blocks beyond the nonzero-polynomial assumption.
\end{theorem}
\begin{proof}
Injectivity removes the inner sums in \eqref{lhs}. Therefore the difference is
the explicit nonnegative sum
\begin{equation}\label{gap}
 \sum_{|\nu|=m}|d_\nu|^2\bigl(W(\nu)-\nu!\bigr).
\end{equation}
The lemma proves the inequality.

For equality, no nonzero $d_\nu$ can have a repeated symbol $y_\alpha$.
Hence $P$ is multiaffine in its $y$ variables. In the polynomial domain,
the degree in $y_\alpha$ of the product defining $P$ is exactly the number
of its linear factors containing that symbol. Thus each $\alpha\in E$ can
occur in at most one $p_i$.

If a monomial of $p_i$ and a monomial of $p_j$, $i\ne j$, overlapped in an
original coordinate, select these and any nonzero monomial from each other
factor. Because symbols belong to unique factors, the resulting $d_\nu$
is their nonzero coefficient product. The lemma then makes its contribution
in \eqref{gap} strictly positive. Thus equality forces disjoint coordinate
supports. Disjoint supports plainly give equality.

Finally, if a nonzero product of linear forms does not involve a coordinate,
none of its linear factors involves that coordinate: degrees in that variable
add in a polynomial domain. Disjoint polynomial supports therefore imply
orthogonal spans for distinct blocks of Gram vectors, so $A_{ij}=0$ for
$i\ne j$. The converse follows directly from block diagonality.
\end{proof}

If a block polynomial is zero, a constituent linear form is zero in the split
case. Then the original Gram matrix has a zero row and both permanents vanish;
the original positive-definite-diagonal-block equality hypothesis excludes this.

\subsection*{Concrete sufficient conditions}
Linear independence over $\mathbb Q$ of the exponent vectors in $E$ implies
injectivity for every $m$. In particular $E=\{k e_1,\ldots,k e_r\}$ works.
For this pure-power case, putting
$P(y)=\sum_{|\nu|=m}d_\nu y^\nu$ as above yields
\[
 \per(A)-\per(G)=\sum_{|\nu|=m}|d_\nu|^2
 \left(\prod_s\frac{(k\nu_s)!}{(k!)^{\nu_s}}-\prod_s\nu_s!\right).
\]

For $k=2$, view monomials $z_s z_t$ as graph edges and $z_s^2$ as loops.
A forest has independent exponent vectors: eliminate a leaf repeatedly.
An odd cycle also has independent exponent vectors, since its vertex equations
force successive edge coefficients to alternate sign and then to vanish.
Thus any union of tree components and components consisting of an odd cycle
with trees attached satisfies the hypothesis. An even cycle has the familiar
alternating-edge relation and does not satisfy it.

For an exact example take
\[
 p_1=z_1(z_2+z_3),\quad p_2=z_2(z_3-z_1),\quad
 p_3=z_3(z_1+i z_2).
\]
The three exponent vectors of $z_1z_2,z_2z_3,z_3z_1$ are independent. Here
\[
 G=\begin{pmatrix}2&-1&1\\-1&2&i\\1&-i&2\end{pmatrix},
 \qquad \per(G)=14,\qquad\left\|p_1p_2p_3\right\|_F^2=88.
\]
The permanent values were independently checked by exact subset dynamic
programming and enumeration of all permutations. The certificate is stored
in \texttt{marcus-collision-free-example.json}.

\section{Why the coefficient argument does not finish the problem}
The monomial map is not expansive on all polynomials in the independent $y$
variables. For example, use the orthonormal quadratic monomials
\[
 \phi_1=x^2/\sqrt2,\quad\phi_2=xy,\quad\phi_3=y^2/\sqrt2.
\]
Then $2\phi_1\phi_3-\phi_2^2=0$, whereas
$\|2y_1y_3-y_2^2\|_F^2=6$. Thus replacing a product of linear forms $P(y)$
by an unrestricted polynomial and claiming norm expansion is false.
This example is not a Marcus counterexample: the displayed quadratic in $y$
is not a product of linear forms.

A second exact barrier comes from orthogonal coordinate groups. Let $B,C$ be
arbitrary PSD $m\times m$ matrices and choose Gram vectors $x_i,y_i$ in
orthogonal spaces. Take block $i$ to consist of $x_i,y_i$. Then
\[
 A\text{ is permutation-similar to }B\oplus C,\quad
 G=B\circ C,
\]
so Marcus in this family becomes
\[
 \per(B)\per(C)\geq\per(B\circ C).
\]
This is Chollet's permanent inequality. The corresponding polynomial support
has the four-cycle relation
$(x_1y_1)(x_2y_2)=(x_1y_2)(x_2y_1)$. Consequently, resolving all collisions
would already handle this substantial additional permanent conjecture; the
absence-of-collisions hypothesis cannot just be omitted.

\section{An exact obstruction to positive cumulants}
For three homogeneous polynomials of the same positive degree, let
\[
 p=\|p_1p_2p_3\|_F^2,\quad q=\per(G),\quad
 H_{ab}=\left\langle\prod_{i\ne a}p_i,\prod_{j\ne b}p_j\right\rangle_F,
 \quad S=\sum_{a,b=1}^3G_{ab}H_{ab}.
\]
Define $D_1=S-3q$ and $D_2=p-S+2q$. Their sum is exactly $p-q$.
If $Z^{(1)},\ldots,Z^{(R)}$ are independent copies and
$U_i=R^{-1/2}\sum_{a=1}^R p_i(Z^{(a)})$, grouping the copy labels gives
\begin{equation}\label{replica}
 \mathbb E\prod_{i=1}^3|U_i|^2
  =q+\frac{D_1}{R}+\frac{D_2}{R^2}.
\end{equation}
To verify this, assignments using one copy contribute $Rp$, those using
two copies contribute $R(R-1)S$, and those using three copies contribute
$R(R-1)(R-2)q$, before division by $R^3$. Unequal holomorphic and
antiholomorphic degrees within any copy give zero by circular symmetry.

Consider the following real integer Gram vectors, paired in row order:
\[
 V=\begin{pmatrix}
 -1&-1&-1\\1&-1&0\\1&1&-1\\1&1&-1\\0&0&1\\0&0&-1
 \end{pmatrix},\qquad A=VV^{\mathsf T}.
\]
Exact arithmetic gives
\[
 G=\begin{pmatrix}6&0&0\\0&18&-2\\0&-2&2\end{pmatrix},\quad
 p=624,\quad q=240,\quad S=1168,
\]
and therefore
\[
 D_1=448,\qquad D_2=-64,\qquad p-q=384.
\]
This disproves termwise nonnegativity of the expansion, not Marcus's
inequality. The PSD condition is exact from the displayed Gram factorization.
The values were checked both by integer subset dynamic programming and
independent permutation enumeration. Full data are in
\texttt{marcus-negative-connected-term.json}.

\section{Numerical work and the remaining gap}
The direct Gram search completed 1,650 random-start descents over eleven
combinations of block size, block count, and rank, up to order 14. It found
no negative difference. Its permanent calculation was checked against direct
permutation sums, and its analytic derivatives against finite differences.
The smallest numerical ratios approached 1 from above. These are local
searches and establish no universal lower bound.

Separate experiments tested the stronger assertion for arbitrary homogeneous
polynomials, and the independent-copy expansion \eqref{replica}. Those tests
also produced no counterexample. Some values below 1 by about $10^{-15}$
are ordinary floating-point roundoff, not violations. In searches for the
ratio $D_2/D_1$, negative local optima around $-0.388$ occurred; no global
bound is claimed, and such a ratio by itself is not a counterexample.

The complete open step is to prove \eqref{lhs}$\geq$\eqref{rhs} in the presence
of colliding exponent sums for every split family $p_i$, and to establish
the equality clause in that generality, or else to give an exact split-family
counterexample. The exact certificate code is ready to verify any strict
numerical witness after rational reconstruction.
The later three-pairs proof handles all collisions for $m=3,k=2$; it does
not remove this remaining gap for general block counts and sizes.

\paragraph{Reproducibility.}
\texttt{marcus-block-search.mjs} contains the main Gram search;
\texttt{marcus-block-exact.mjs} contains independent exact checks;
\texttt{marcus-chaos-search.mjs} explores arbitrary homogeneous polynomials;
\texttt{marcus-replica-search.mjs} explores the cumulant terms.
\texttt{marcus-search-summary.json} indexes the direct-search results.
Lean formalization has not been attempted, as requested. The goal remains active.
\end{document}
